\chapter*{\centerline{Abstract}}
\markboth{\MakeUppercase{Abstract}}{}
\iftoggle{fulltoc}{
  \addcontentsline{toc}{chapter}{Abstract}
}{}

Biomimetic robots whose legs reproduce human muscle geometry and joint torque would allow researchers to run neuromechanical experiments that are impractical or unethical in living subjects, but designing them requires tools that predict artificial-muscle force, joint torque, and neural control before hardware exists. This dissertation develops and validates those tools for the humanoid robot under construction at Portland State University, a robot actuated from the trunk down by Festo braided pneumatic actuators (BPAs) and controlled by a synthetic nervous system.

The isometric force of BPAs with \qtylist{10;20}{\mm} uninflated diameters was characterized across resting lengths, contractions, and pressures on a custom test jig. The characterization produced two models: a maximum-force law at \qty{620}{\kPa} whose resting-length dependence was previously unreported and is predicted with the wrong trend by the manufacturer's tool, and a three-coefficient normalized force surface (adjusted $R^2 > 0.99$) that corrects the large over-predictions of existing models at short resting lengths.

The actuator model was then carried to the joint through isometric knee torque experiments on a pinned knee and on a biomimetic four-bar knee with a migrating instantaneous center of rotation. A hybrid torque calculation method isolated three loss mechanisms: constant length offsets, series compliance in brackets and tendons, and loss of usable actuator length where BPAs wrap the joint. Correction terms identified by multiobjective optimization improved prediction on every validation test, with per-test error falling from 1.0--3.6 to 0.6--1.9~N{\,}m across flexor and extensor configurations. With actuator placements redesigned under the corrected model, the biomimetic knee met or exceeded the human isometric torque envelope over \textbf{[FLAG-BEN: \% of RoM]} of the range of motion (RMSE \textbf{[FLAG-BEN: value]}~N{\,}m against the human target).

Finally, the validated plant was connected to biologically grounded control. The laboratory's modified OpenSim leg model was converted to MuJoCo, and a synthetic nervous system implementing a two-layer central pattern generator, with per-leg rhythm generation and pattern formation, left--right coordination, and Ia, Ib, II, and heel/toe contact feedback, was verified against published tonic-stimulation and deletion phenomena. With these tools, researchers can now size artificial muscles and their routing to human torque specifications, validate the resulting joints against measurement, and run lesion and stimulation experiments on the simulated humanoid leg in preparation for a treadmill robot with embedded real-time neural control.
